\chapter{Se servir de l'atelier}

Les chapitres précédents produisaient du code. Celui-ci apprend à s'en servir —
et surtout à \textbf{lire un résultat sans se tromper}, ce qui est plus difficile
que de l'obtenir.

\section{Les sept panneaux}

Tous sont dockables : on les déplace, on les empile en onglets, on les ferme.

\begin{center}
\begin{tabular}{@{}p{2.4cm}p{11.4cm}@{}}
\toprule
\textbf{Panneau} & \textbf{Ce qu'il fait} \\
\midrule
Plateau & l'état, le picking, les coups légaux, la menace, le dernier coup, le score \\
Règles & la grille $+$ tous les paramètres du moteur, auto-générés \\
Joueurs & qui tient chaque siège, à quelle force, avec quel budget \\
Modules & les \texttt{.cpp} détectés, leur statut, \textbf{la sortie du compilateur} \\
Journal & la suite des coups, leur empreinte, le curseur de rejeu \\
Métriques & la campagne IA contre IA, et ce qu'elle répond \\
Sortie & \textbf{ce que votre code écrit} quand il tourne (\texttt{NKC\_LOG\_*}) \\
\bottomrule
\end{tabular}
\end{center}

\begin{nkpiege}[Un panneau qui n'est pas l'onglet actif ne se dessine pas]
C'est le fonctionnement normal du docking. Conséquence : ouvrez le panneau
\emph{Métriques} pour voir la progression d'une campagne. La campagne, elle,
tourne quand même — elle vit sur ses propres threads, pas dans l'interface.
\end{nkpiege}

\section{Lire le plateau}

\begin{itemize}
  \item \textbf{disques colorés} : les totems. Le \emph{niveau} se lit à la
        \emph{taille} et au liseré, jamais à la seule couleur — celle-ci porte
        déjà le propriétaire ;
  \item \textbf{anneaux verts} : les destinations légales du totem sélectionné,
        avec une étiquette qui dit le \emph{genre} du coup --- \texttt{D}
        (dupliquer), \texttt{F3} (fusionner trois cases), \texttt{P1}
        (pouvoir n\textsuperscript{o}1) ;
  \item \textbf{anneaux orange reliés par un trait} : les cases qu'une
        \textbf{fusion} va consommer, et au bout du trait la case du
        \textbf{résultat}, en anneau plein plus épais ;
  \item \textbf{« $\times$2 » au-dessus d'une case} : plusieurs coups
        différents visent cette case. Cliquez : l'atelier demandera lequel ;
  \item \textbf{anneaux rouges} : les totems ennemis que le coup \emph{survolé}
        retournerait ;
  \item \textbf{anneau orange qui pulse} : le dernier coup joué, une seconde
        environ ;
  \item \textbf{« CASCADE $\times$N »} au centre : deux totems ou plus retournés
        d'un coup.
\end{itemize}

Les anneaux rouges sont la lecture centrale du jeu : « quelle surface de contact
suis-je en train d'offrir ? ». Ils ne sont pas déduits des règles — l'atelier
simule le coup et lit les événements. Ils restent donc justes même quand vous
changez la règle de transformation.

\section{Jouer un coup à la souris}

Dupliquer, fusionner, lancer un pouvoir : \textbf{un clic sur une case de
départ, un clic sur la case d'arrivée}. Le geste ne change jamais ; ce qui
change, c'est quelle case est le départ et quelle case est l'arrivée --- et
pour une fusion, cela ne se devine pas.

Tout est au chapitre suivant, \textbf{«~Jouer un coup à la souris~»} : où
cliquer pour jouer une \textbf{fusion}, comment jouer un \textbf{pouvoir} et
en choisir un quand plusieurs visent la même case, et ce qui existe réellement
pour les \textbf{artefacts}.

\subsection{La barre d'état, et les boutons}

La barre dit à tout instant ce que l'atelier attend. \textbf{Relisez-la avant de
chercher un bug.}

\begin{center}
\begin{tabular}{@{}p{4.2cm}p{9.6cm}@{}}
\toprule
\textbf{Bouton} & \textbf{Effet} \\
\midrule
Nouvelle partie & rejoue depuis le début, mêmes réglages, même graine \\
Lecture / Pause & la partie avance-t-elle toute seule ? \\
Pas à pas & un seul coup d'IA, même en pause \\
IA vs IA & tous les sièges à l'IA $+$ lecture : la configuration de mesure \\
Siège : Humain / IA & bascule le siège \textbf{au trait} \\
Voisinage & trace les liens d'adjacence de la case survolée (mise au point) \\
\bottomrule
\end{tabular}
\end{center}

\section{Le panneau Sortie : vos propres traces}

Le panneau \textbf{Modules} montre ce que dit le \emph{compilateur} ; celui-ci
montre ce que dit votre \emph{code} une fois qu'il tourne. Les deux sont
nécessaires, et les mélanger rendrait les deux illisibles.

\begin{itemize}
  \item \textbf{Niveau minimal} filtre : passez à \texttt{WARN} quand vous
        cherchez une anomalie dans un flot de \texttt{INFO} ;
  \item les lignes identiques \textbf{consécutives} sont fusionnées en
        \texttt{(xN)} — une boucle qui répète la même trace mille fois reste
        lisible ;
  \item le plus récent est \textbf{en haut}, donc toujours visible sans défiler ;
  \item le tampon est borné à \textbf{4096 lignes}, et le panneau \textbf{affiche
        ce qu'il a perdu}. Un journal qui perd des lignes en silence fait chercher
        un bug qui n'existe pas.
\end{itemize}

\begin{nkpiege}[Les modules internes n'y écrivent pas]
\texttt{ConquerorRulesV2 (interne)} et \texttt{AIRef (interne)} sont compilés
\emph{dans} l'application : ils partagent son journal, et leurs traces vont dans
la console, pas ici. Le panneau Sortie est là pour \textbf{votre} code.
\end{nkpiege}

\section{Reproduire un bug : le journal}

C'est la fonction la plus utile de l'atelier, et la moins spectaculaire.

Le panneau \emph{Journal} liste les coups avec, à droite de chacun,
\textbf{l'empreinte de l'état après ce coup}. Le bouton \textbf{Copier} met dans
le presse-papiers une trace rejouable : graine, nombre de joueurs, un coup par
ligne, empreinte finale.

\begin{nkretenir}[Comment on s'en sert vraiment]
Votre moteur produit un résultat différent sur Linux et sur Windows. Vous copiez
le journal des deux côtés, vous comparez les colonnes d'empreintes, et vous
trouvez \textbf{la première ligne qui diffère}. Ce coup-là, et pas un autre,
contient votre non-déterminisme.

Sans les empreintes intermédiaires, vous sauriez seulement que les états finaux
diffèrent — c'est-à-dire presque rien.
\end{nkretenir}

Le bouton \textbf{Copier} produit une trace \textbf{auto-descriptive} :

\begin{nkcode}{Entete d'une trace copiee}
# ConquerorLab — trace rejouable
# graine        20260807
# joueurs       2
# regles        ConquerorRulesV2 (interne)
# siege 0       Humain
# siege 1       AIRef (1.0.0)
# plateau       hexagone_6x7.json
#
# joueur action de_q de_r vers_q vers_r empreinte
\end{nkcode}

La graine et les coups ne suffisaient pas : rejouée avec un autre moteur ou une
autre IA, la même liste donne un autre résultat, et on cherche un bug là où il
n'y en a pas. Vous êtes deux à travailler en parallèle — sans cet entête, un
rapport ne dit même pas de qui vient le code qu'il incrimine.

Le curseur de rejeu \textbf{reconstruit} la position (\texttt{Setup} $+$ N coups)
au lieu de « défaire » le dernier coup : le contrat n'impose au moteur aucune
opération inverse, et une annulation approximative donnerait un état qui n'a
jamais existé.

\section{Deux garde-fous contre les erreurs silencieuses}

Ce sont les plus dangereuses : rien ne plante, les chiffres sortent, et ils sont
faux.

\subsection{L'incohérence voisinage / coups légaux}

Un moteur déclare son voisinage (\texttt{GetNeighbors}) et génère ses coups. Rien
n'oblige les deux à s'accorder. Quand ils divergent, \textbf{rien ne se plaint} —
mais une IA qui compte « combien d'ennemis touche cette case » se trompe, et joue
simplement moins bien. Une semaine perdue à chercher dans l'évaluation un bug qui
est dans la géométrie.

L'atelier vérifie donc, à chaque nouvelle partie, qu'un coup DUPLIQUER va bien
vers une case que \texttt{GetNeighbors} déclare voisine. Sinon, bandeau rouge en
tête du panneau \emph{Modules} :

\begin{nkcode}{Ce que vous lisez si les deux se contredisent}
INCOHERENCE : 3 coup(s) DUPLIQUER visent une case que GetNeighbors ne
declare pas voisine — p.ex. (2,1) -> (4,1). Ton generateur de coups et ton
voisinage ne disent pas la meme chose ; toute IA qui evalue l'adjacence se
trompera.
\end{nkcode}

Il ne dit pas lequel des deux a tort — il signale la contradiction, ce qui suffit
à la faire chercher.

\subsection{La signature de configuration}

\begin{nkcode}{Deux chiffres, et aucun moyen de les comparer}
campagne du matin  : 62 % pour le camp 0
campagne du soir   : 51 % pour le camp 0
\end{nkcode}

Que s'est-il passé ? A1 a changé un paramètre ? A2 son évaluation ? Les deux ?
Sans réponse, ces deux nombres ne veulent \textbf{rien} dire — et on ne peut même
pas savoir qu'ils ne veulent rien dire.

Le panneau \emph{Métriques} affiche donc, \textbf{au-dessus} des chiffres, tout
ce qui les a produits : moteur, IA de chaque siège, paliers, budget, plateau,
valeurs de \textbf{tous} les paramètres, plus un identifiant court. Et il compare
avec la campagne précédente :

\begin{nkcode}{Le bandeau qui evite une conclusion fausse}
Depuis la campagne precedente, les PARAMETRES, les IA ont change.
Les deux resultats ne se comparent pas.
\end{nkcode}

\begin{nkretenir}[Pourquoi ça marche là où la discipline échoue]
« Une variable à la fois » est la bonne règle. Mais une discipline qu'on ne peut
pas \textbf{vérifier après coup} n'en est pas une : personne ne se souvient, une
semaine plus tard, de ce qui avait bougé.

On n'empêche personne de changer deux choses ; on rend impossible de ne pas s'en
apercevoir. C'est la seule forme de garde-fou qui tienne entre deux personnes.
\end{nkretenir}

\section{La campagne : ce qu'elle mesure, et comment}

Panneau \emph{Métriques}. Réglez le nombre de parties, les threads, le budget, les
IA de chaque camp, puis \textbf{Lancer la campagne}.

\subsection{L'inversion des côtés}

\begin{nkcode}{src/ConquerorLab/NkcBatch.h — en-tete}
2. LES COTES SONT INVERSEES une partie sur deux. Sans cela, tout ecart
   mesure melange « cette IA est meilleure » et « le premier joueur est
   avantage » -- et REGLES 15 pose justement l'avantage au premier joueur
   comme LA question du palier 0. On mesure les deux separement.
\end{nkcode}

Concrètement : les parties $2n$ et $2n+1$ partagent la \textbf{même graine} et
n'échangent que les sièges. C'est cet appariement qui isole l'avantage de
position. D'où deux tuiles distinctes :

\begin{itemize}
  \item \textbf{Camp 0 / Camp 1} : quelle \emph{stratégie} gagne, sièges
        confondus ;
  \item \textbf{Siège 1 gagne} : quel \emph{siège} gagne, stratégies confondues.
        C'est l'avantage au premier joueur.
\end{itemize}

\begin{nkpiege}[L'inversion ne s'applique qu'à deux joueurs]
\begin{quote}\small
À 2 joueurs l'inversion échange 0 et 1 ; au-delà on ne l'applique pas — une
permutation cyclique à 3--4 joueurs ne s'apparie plus deux à deux.
\hfill\textsc{NkcBatch.h, SeatToConfig}
\end{quote}
À 3 ou 4 joueurs, les winrates par camp \textbf{mélangent} effet de stratégie et
effet de siège. Tenez-en compte avant de conclure quoi que ce soit.
\end{nkpiege}

\subsection{Les tuiles rouges}

\textbf{Coupées (\texttt{max\_tours})} doit valoir \textbf{zéro}. Dès que la tuile
est rouge, des parties ne se terminent pas d'elles-mêmes — et « un taux non nul
signale une pathologie des règles, pas un réglage à monter ».

\textbf{Coups illégaux} doit valoir zéro aussi. Un chiffre non nul veut dire
qu'une IA a proposé un coup que le moteur refuse — presque toujours un
\texttt{memset} manquant, parfois une IA qui construit un coup au lieu de le
choisir.

\subsection{Les deux histogrammes}

\textbf{Usage des actions} — cinq barres : \emph{aucune}, \emph{dupliquer},
\emph{fusionner}, \emph{pouvoir}, \emph{passer}. La barre \emph{aucune} doit
rester à zéro : une valeur non nulle serait un bug du générateur de coups, pas un
détail à masquer.

Une action jamais jouée est une action à supprimer ou à rendre attractive. C'est
ce raisonnement qui a tué l'action SAUTER :

\begin{quote}\small
L'action SAUTER n'est jamais jouée : $+0$ totem, une seule cible, conquête
impossible — dominée par DUPLIQUER dans toute position. \textbf{Supprimée} du jeu
de base. Revient comme \emph{pouvoir}, où elle a un coût.
\hfill\textsc{regles\_completes\_v2.md:57}
\end{quote}

\textbf{Durée des parties} — une barre par tranche de dix coups. Une distribution
très étalée, ou un pic contre la dernière barre, signale des parties qui traînent.

\section{Une méthode de mesure qui tient}

Quatre règles, apprises à leurs dépens par d'autres.

\begin{enumerate}
  \item \textbf{Une variable à la fois.} Changez \texttt{portee\_duplication},
        mesurez, notez. Puis remettez-la et changez autre chose. Deux variables
        ensemble donnent un résultat que rien ne permet d'attribuer.
  \item \textbf{Assez de parties.} Sur 20 parties, un écart de 10 points n'est que
        du bruit. Comptez 200 pour une tendance, 1000 pour une conclusion, 10\,000
        pour une décision de conception.
  \item \textbf{Notez la graine et les réglages avec le chiffre.} Un winrate sans
        son contexte n'est pas une mesure, c'est une anecdote. La barre du plateau
        affiche la graine en permanence pour cette raison.
  \item \textbf{Le résultat qui ne bouge pas est un résultat.} Si doubler la
        portée ne change pas le winrate, vous venez d'apprendre que ce paramètre
        n'est pas un levier. C'est une information, pas un échec.
\end{enumerate}

\begin{nkpiege}[Le réglage modifié en cours de partie]
Le panneau \emph{Règles} affiche un bandeau orange dès que vous touchez un
paramètre : « Réglage modifié — relance la partie pour mesurer ». Il est là parce
que les coups \textbf{déjà joués} ne sont pas rejoués avec la nouvelle règle. Une
partie à cheval sur deux réglages ne mesure rien.
\end{nkpiege}

\section{Quand ça ne marche pas}

\begin{center}
\begin{tabular}{@{}p{6.0cm}p{7.8cm}@{}}
\toprule
\textbf{Symptôme} & \textbf{Regardez d'abord} \\
\midrule
rien ne bouge & \textbf{la barre d'état du plateau} — elle dit ce qu'on attend \\
mon module n'apparaît pas & panneau \emph{Modules} : détecté ? compilé ? bon chemin ? \\
« Symbole introuvable » & les deux \texttt{NKC\_MODULE\_EXPORT} en fin de fichier \\
« ABI 2, attendue 3 » & vous compilez contre de vieux en-têtes \\
« coup ILLÉGAL » & le \texttt{memset} de votre coup \\
l'interface se fige une seconde & c'est une compilation, pas un plantage \\
« le moteur a REFUSÉ ce plateau » & \texttt{cells} vide, JSON malformé, ou moteur sans \texttt{LoadBoardJson} \\
la campagne refuse de démarrer & chaque siège doit être tenu par une IA \emph{chargée} \\
\bottomrule
\end{tabular}
\end{center}

\section{Le banc d'essai, hors interface}

Avant de croire à quoi que ce soit, faites tourner le banc d'essai. Il vérifie le
contrat de bout en bout, sans interface :

\begin{nkcode}{Applications/ConquerorLab/README.md — section « Reproduire »}
CLANG="C:/msys64/ucrt64/bin/clang++.exe"
R="d:/Projets/2026/Nkentseu/Nkentseu"
INC="-I$R/Applications/ConquerorLab/include \
     -I$R/Kernel/Foundation/NKCore/src \
     -I$R/Kernel/Foundation/NKPlatform/src"

$CLANG -shared -std=c++17 -O2 -fPIC $INC -o ConquerorRulesV2.dll \
       $R/Applications/ConquerorLab/modules/rules/ConquerorRulesV2.cpp
$CLANG -shared -std=c++17 -O2 -fPIC $INC -o ConquerorAIRef.dll \
       $R/Applications/ConquerorLab/modules/ai/ConquerorAIRef.cpp
$CLANG -std=c++17 -O1 $INC -o abitest.exe \
       $R/Applications/ConquerorLab/tests/NkcAbiHarness.cpp
./abitest.exe
\end{nkcode}

Seize vérifications, dont les trois qui comptent : l'IA ne produit jamais de coup
illégal, la partie se termine d'elle-même, et le rejeu redonne exactement la même
empreinte.

Remplacez les deux chemins de modules par les vôtres : le banc d'essai teste
\textbf{votre} module aussi bien que celui de référence.

\begin{nkexo}[1 — Le premier chiffre]
Lancez 200 parties, IA de référence \emph{Normal} contre \emph{Facile}, sur le
plateau par défaut. Notez : winrate par camp, victoires du siège 1, taux de
coupées, durée moyenne. Recommencez avec une autre graine. Les chiffres
bougent-ils ? De combien ? Vous venez de mesurer votre bruit de fond — donc le
seuil en dessous duquel un écart ne veut rien dire.
\end{nkexo}

\begin{nkexo}[2 — L'avantage au premier joueur]
Même IA des deux côtés, 1000 parties. Que vaut « Siège 1 gagne » ? Est-ce
significatif au regard du bruit mesuré à l'exercice 1 ? C'est \textbf{la} question
du palier 0 : écrivez votre réponse en une phrase, avec le chiffre.
\end{nkexo}

\begin{nkexo}[3 — Le paramètre qui ne sert à rien]
Prenez un paramètre du moteur de référence et faites-le varier sur trois valeurs,
300 parties chacune. Tracez le winrate. Est-ce un levier ou une décoration ? Sur
quel critère tranchez-vous ?
\end{nkexo}

\begin{nkexo}[4 — Reproduire une divergence]
Introduisez volontairement du non-déterminisme dans votre moteur : par exemple,
faites dépendre l'ordre des coups de l'adresse d'un pointeur. Jouez une partie,
copiez le journal, relancez avec la même graine, comparez les empreintes. À quel
coup la divergence apparaît-elle ? Retirez ensuite le défaut.
\end{nkexo}
